% Minimal retained preamble for S3 SIGMA 2022 article 027 cards taken
% from sources/sigma-2022-027-twistor-theory-of-dancing-paths/
% pinned-original-source.tex (doi:10.3842/SIGMA.2022.027, SIGMA 18
% (2022), CC BY 4.0).
%
% Omitted and why: the publisher documentclass{sigma} block (source
% lines 1-7, including the ifpdf probe) and the frontmatter commands -
% the sigma class is unavailable in the toolchain and the assembler
% supplies the article class. The author equation numbering (source
% line 9), the author theorem declarations (source lines 11-19,
% including the starred Theorem* and the definition-style group with
% its brace pair on lines 16 and 19) and the author macro subset used
% by the delivered spans (source lines 41-108: \CP \C \Z \PP \RP \tL
% \tP \R \Rho \p \be \u \z \I \ee \ep \spp \tr \eps, the line-90
% redefinition of \Z to \mathcal{Z}, \U \cI \st \uh \zh) are copied
% verbatim below in source order. Omitted from the macro block: the
% margin-note macros (lines 21-40: \mnotex \mnote \md \tm \hook) and
% unused shorthands (\projb \bilb \bX \mf \bw \L \d \bv \a \b \r \l
% \ov \eqq \bea \eea \hg \hw \hnabla \hk \hrho \hL \hpsi \he \hS \T \E
% \ush \zsh \rk \ms \ent \bs \bn \bnbf \n) - none occurs in the
% delivered spans; the page-style \allowdisplaybreaks (line 111) is
% layout-only. The qed command is already provided by the
% amsthm/assembler baseline; the author's verbatim \renewcommand{\qed}{}
% at the end of one delivered proof compiles unchanged.
\usepackage{amsmath,amssymb,amsthm}
\numberwithin{equation}{section}

\newtheorem{Theorem}{Theorem}[section]
\newtheorem*{Theorem*}{Theorem}
\newtheorem{Corollary}[Theorem]{Corollary}
\newtheorem{Lemma}[Theorem]{Lemma}
\newtheorem{Proposition}[Theorem]{Proposition}
{ \theoremstyle{definition}
\newtheorem{Definition}[Theorem]{Definition}
\newtheorem{Example}[Theorem]{Example}
\newtheorem{Remark}[Theorem]{Remark} }

\newcommand{\CP}{\mathbb{CP}}
\newcommand{\C}{\mathbb{C}}
\newcommand{\Z}{\mathbb{Z}}
\newcommand{\PP}{\mathbb{P}}
\newcommand{\RP}{\mathbb{RP}}
\newcommand{\tL}{\widetilde{L}}
\newcommand{\tP}{\widetilde{P}}
\newcommand{\R}{\mathbb{R}}
\newcommand{\Rho}{\mathrm{P}}
\def\p{\partial}
\def\be{\begin{equation}}
\def\u{\mathfrak{u}}
\def\z{\mathfrak{z}}
\def\I{\mathcal{I}}
\def\ee{\end{equation}}
\def\ep{{\varepsilon}}
\newcommand{\spp}{\mathbb{S}}
\newcommand{\tr}{\operatorname{tr}}
\newcommand{\eps}{\varepsilon}
 \renewcommand{\Z}{\mathcal{Z}}
 \newcommand{\U}{\mathcal{U}}
\newcommand{\cI}{{\mathcal I}}
\newcommand{\st}{\mid }
\newcommand{\uh}{\widehat{u}}
\newcommand{\zh}{\widehat{z}}
